\subsection{Hellinger and statistical distance}

The squared Hellinger distance between distributions $D_1$ and $D_2$ over a discrete probability space
${\cal A}$ is
\[
H^2(D_1,D_2) := \frac{1}{2} \sum_{a\in {\cal A}}\left(\sqrt{D_1(a)}
  - \sqrt{D_2(a)}\right)^2 = 1 - \sum_{a\in {\cal A}}\sqrt{D_1(a)
  D_2(a)}\,.
\]
It is clear that $1-H^2(\cdot , \cdot)$ is multiplicative for product distributions $D_1^k, D_2^k$ on space
${\cal A}^k$, \ie,
  $$ 1-H^2(D_1^k, D_2^k) =  (1-H^2(D_1,D_2))^k\,.   $$
The statistical distance between $D_1$ and $D_2$ is:
$$\Delta(D_1,D_2) :=  \frac{1}{2} \sum_{a\in {\cal A}}\left|D_1(a) - D_2(a)\right|\,.$$
We have the standard inequality:
\begin{lemma}[\cite{Pol}]\label{p:Hellinger} 
$\displaystyle H^2(D_1,D_2) \leq \Delta(D_1,D_2) \leq \sqrt{2} \cdot H(D_1,D_2)\,.$
\end{lemma}









\section{The Outer PCP}\label{sec:outer}
In this section, we describe our Outer PCP\@. For the ease of exposition, we do it in three stages,
starting with the standard \emph{Variable Versus Equation} Game. We also prove the \emph{sub-code covering property}.
Let $(X, \Phi)$ be an instance of the 3LIN problem over $\F_q$ given by \expref{Theorem}{thm:Has97}
or \expref{Theorem}{thm:kp}.


\subsection{The Variable Versus Equation Game}

Consider the 2P1R Game where the verifier picks a random equation $E \in \Phi$ and then picks one of the three
variables $x \in E$. The first prover is sent the question $E$ (\ie, the three variables appearing in $E$)
and the second prover is sent the question $x$. The provers answer with the $\F_q$-values of all the variables
they receive.
The verifier accepts if and only if the two provers agree on the variable $x$ and moreover
the values given by the first prover
      satisfy the equation $E$. It is well-known and easy to check that
if $\OPT(X,\Phi) = 1-\eps$, then the value of the game is $1- {\eps}/{3}$.

\subsection{The Basic 2P1R Game}

We now slightly modify the Variable Versus Equation Game and call it the \emph{Basic 2P1R Game}
(think of $\beta$ as small):


\begin{itemize}
\item The verifier picks an equation $E \in \Phi$ at random. Let $x, y, z \in X$ be the variables in $E$.
\item The first prover is sent the equation $E$.
\item The second prover is sent the equation $E$ with probability $1-\beta$ and one of the three
variables $x,y, z$ with probability ${\beta}/{3}$ each.
\item The provers answer with the values of all the variables they receive.
\item The verifier accepts if and only if the two provers agree on the
      values of the variables and moreover the values given by the first prover
      satisfy the equation.
\end{itemize}

Assume that $\OPT(X,\Phi) = 1-\eps$.  The value of the game is at least $1-\eps$ since the provers
may stick to a $(1-\eps)$-satisfying assignment to $(X,\Phi)$ and in that case, the verifier may reject only
when the equation $E$ is not satisfied.

On the other hand, the value of the game is at most
$1-\Omega(\eps \beta)$ since with probability $\beta$, the second prover receives exactly one
variable and the variable versus equation game is played.

\subsection{The Final 2P1R Game (Outer PCP)}

The Outer PCP is now obtained as the $k$-wise repetition applied to the Basic 2P1R Game. Specifically:

\begin{itemize}
\item The verifier picks equations $E_1,\ldots,E_k \in \Phi$ at random. Let $V$ denote the set of $3k$ variables appearing in these equations.
\item The question $V$ is sent to the first prover.
\item Let $U \subseteq V$ be chosen by including independently for each $1 \leq i \leq k$, all three variables in equation $E_i$ with probability $1-\beta$ and one of the three variables in equation
    $E_i$ with probability ${\beta}/{3}$ each. The question $U$ is sent to the second prover.
\item The provers answer with the values of all the variables they receive.
\item The verifier accepts if and only if the two provers agree on the
      values of the variables in $U$ and moreover the values given by the first prover
      satisfy all the $k$ equations $E_1, \ldots, E_k$.
\end{itemize}

Assume again that $\OPT(X,\Phi) = 1-\eps$.  The value of the game is at least $1-\eps k $
since the provers
may stick to a $(1-\eps)$-satisfying assignment to $(X,\Phi)$ and in that case, the verifier may reject only
when at least one equation $E_i$ is not satisfied.

On the other hand, we observe that the value of the game is at most $(1-\Omega(\eps^2))^{\Omega(\beta k)}$.
As noted before, the value of the Basic Game is $1-\Omega(\eps \beta)$ and if we apply the parallel repetition
theorem (\expref{Theorem}{thm:pr}) directly, we end up with an upper bound of
$(1-\Omega(\eps^2 \beta^2))^{\Omega(k)}$. This bound is not good enough for us and
we get the better bound as follows.
Call a coordinate useful if the second prover receives a single variable, \ie,
the standard Variable Versus Equation Game is played. The expected number of useful
coordinates is $\beta k$. Hence the probability that less than $\beta k/2$
coordinates are useful is at most $2^{-\Omega(\beta k)}$ and may be ignored in comparison to the desired bound.
The repeated game restricted to the question pairs where at least $\beta k/2$
coordinates are useful can be thought of as a convex combination of sub-games,
each sub-game being the standard Variable Versus Equation Game repeated at least
$\beta k/2$ times. Each such sub-game has value at most $(1-\Omega(\eps^2))^{\Omega(\beta k)}$
and hence
so does the overall repeated game.

We note that the $k$ equations chosen by the verifier in the first step are
disjoint (\ie, do not share variables)
with high probability and we make this assumption in all that follows.


\subsection{Sub-code covering property}

Fix a question $(E_1, \ldots, E_k)$ to the first prover in the Outer PCP and
let $V$ denote the set of variables $\{x_1,\ldots,x_{3k}\}$ in these equations.
The question to the second prover is a subset $U \subseteq V$ where independently
for each $1 \leq i \leq k$,
all three variables
$x_{3(i-1)+1}, x_{3(i-1)+2}, x_{3i}$ are included in $U$ with probability $1-\beta$ and
exactly one of the three variables is included with probability ${\beta}/{3}$ each.
Consider two
distributions on $\F_q^V = \F_q^{3k}$:
\begin{itemize}
\item Distribution ${\cal D}$ is uniform on $\F_q^{V}$.
\item Distribution ${\cal D'}$
%
%
%
is obtained by picking
a random question $U \subseteq V$ to the second prover (as described above),
picking a string
in $\F_q^U$ uniformly at random, and then ``pulling up'' this string to
$\F_q^V$ by inserting $0$ at positions in $V \setminus U$.
\end{itemize}

\begin{lemma}\label{lemma:covering}  The statistical distance
$\Delta( {\cal D}, {\cal D'} )$ is upper bounded by $O(\sqrt{k}\cdot q\beta)$.
\end{lemma}
\begin{proof} We will bound the squared Hellinger distance on one coordinate, then
use the multiplicativity of the squared Hellinger distance for product distribution,
and then upper bound the statistical distance in terms of the Hellinger distance.

Let ${\cal Q}$ denote the uniform distribution on $\F_q$ and $({\cal Q}, {\cal Q}, {\cal Q})$
denote its three independent copies.
On any single coordinate $1 \leq j \leq k$,
note that ${\cal D}_j$ is same as the distribution $({\cal Q}, {\cal Q}, {\cal Q})$ whereas
${\cal D}'_j$ can be written as:
 $$ {\cal D}'_j = (1-\beta)\cdot ({\cal Q}, {\cal Q}, {\cal Q})+ \frac{\beta}{3} \cdot
 ({\cal Q}, 0, 0) + \frac{\beta}{3}\cdot (0, {\cal Q}, 0)+ \frac{\beta}{3} \cdot
 (0,0, {\cal Q})\,. $$
The total probability mass attached by ${\cal D}_j$ to triples in $\F_q^3$ where the
number of non-zero entries is three, two, one, and zero respectively is:
$$ p_3=\frac{(q-1)^3}{q^3}, \ p_2= \frac{3(q-1)^2}{q^3}, \ p_1=\frac{3(q-1)}{q^3}, \ p_0=\frac{1}{q^3}\,. $$
Similarly, the total probability mass attached by ${\cal D}'_j$ to triples in $\F_q^3$ where the
number of non-zero entries is three, two, one, and zero respectively is:
$$ p_3'=(1-\beta)\frac{(q-1)^3}{q^3}, \ p_2'=(1-\beta)\frac{3(q-1)^2}{q^3}, \ p_1'=(1-\beta)\frac{3(q-1)}{q^3}+
\beta \frac{q-1}{q}, \ p_0'=(1-\beta)\frac{1}{q^3}+ \beta \frac{1}{q}\,. $$
Hence
\begin{align*}
  1 - H^2({\cal D}_j, {\cal D}'_j )\
  &=\  \sqrt{p_3 p_3'} + \sqrt{p_2 p_2'} + \sqrt{p_1 p_1'} + \sqrt{p_0 p_0'} \\
  &=\  \frac{(q-1)^3}{q^3} \sqrt{1-\beta} + \frac{3(q-1)^2}{q^3}
  \sqrt{1-\beta}\\
  &\phantom{=\ \frac{(q-1)^3}{q^3} \sqrt{1-\beta}\ } +
  \frac{3(q-1)}{q^3} \sqrt{1+ \left(\frac{q^2}{3}-1\right)\beta} + \frac{1}{q^3} \sqrt{1+ (q^2-1)\beta} \\
  &\geq\ 1- O(\beta^2 q^2)\,,
\end{align*}
where all the square roots are estimated using
\[
\sqrt{1+x} \geq 1+\frac{x}{2}-x^2
\]
 for $x \in
[-{1}/{2},{1}/{2}]$
and $q^2 \beta \ll {1}/{2}$ (as will be the case). Thus the expression is
lower bounded by $a_0 + a_1 \beta + a_2 \beta^2$ where it is easily checked that
$a_0 =1 $ and $|a_2| = O(\beta^2 q^2)$.
The main point is that the coefficient $a_1$ vanishes. Indeed, $a_1$ equals
$$ \frac{1}{2} \left( - \frac{(q-1)^3}{q^3} -  \frac{3(q-1)^2}{q^3}
+  \frac{3(q-1)}{q^3} \cdot \left(\frac{q^2}{3}-1\right) +  \frac{1}{q^3} \cdot
  (q^2-1)  \right)   = 0\,.  $$
%
%
%
%
By multiplicativity of the squared Hellinger distance, we have
 $$ 1- H^2({\cal D}, {\cal D}') \ = \ (1- H^2({\cal D}_j, {\cal D}'_j))^{k}
 \ \geq \
 (1-O(\beta^2 q^2))^{k} \ \geq \ 1-O(\beta^2 q^2k)\,.$$
Finally $\Delta({\cal D}, {\cal D}') \leq  \sqrt{2} \cdot
 H({\cal D}, {\cal D}') \leq O(\sqrt{k} \cdot q\beta)$.
\end{proof}



\subsection{The choice of parameters}\label{sec:param}

Let $\OPT(X,\Phi)=1-\eps$ where $(X,\Phi)$ is a NO instance of 3LIN given by \expref{Theorem}{thm:Has97} or
\expref{Theorem}{thm:kp}.
\begin{itemize}
\item In \expref{Theorem}{thm:Has97}, $\OPT(X,\Phi)$ is close to ${1}/{q}$ and hence
$\eps = \Omega(1)$.
\item In \expref{Theorem}{thm:kp},
$\OPT(X,\Phi) \leq 1-\Omega({1}/{\log^3 N})$ where $N$ is the
size of the 3LIN instance and hence $\eps = \eps(N) =  \Omega({1}/{\log^3 N})$.
\end{itemize}
The parameters will be chosen so that for a large enough constant $C$,
\begin{itemize}
\item The soundness of the Outer PCP, which is at most $(1-\Omega(\eps^2))^{-\Omega(\beta k)}$,
      is at most ${1}/{q^{C}}$.
\item $\displaystyle \Delta({\cal D}, {\cal D}') \leq \frac{1}{C q^2}$.
\end{itemize}
Using \expref{Lemma}{lemma:covering},
it suffices to choose 
\[
k = \frac{C_*^3  q^6 \log^2 q}{\eps^4} \quad \text{and}\quad
\beta = \frac{\eps^2}{C_*^2 \cdot q^6 \cdot \log q}
\]
for a large enough constant $C_*$.







%
%
%
%
%
%
%
%
%
%
%
%
%
%



\section{The Inner PCP}

In this section, we describe our Inner PCP\@. We first analyze a test that we call Gowers Test,
which is then used to analyze the actual Inner PCP presented in \expref{Section}{sec:point-subspace}. We begin with
some notation and a simple lemma.

Let $\omega := e^{2\pi i/q}$ be the complex $q^{th}$ root of unity and $\Omega := \{1, \omega,
\ldots, \omega^{q-1} \}$. For $f, g : \F_q^m \mapsto \Omega$, let $\AGR(f,g)$ denote the agreement
between the two functions, \ie, the fraction of points on which they agree.
Let $f_i$ denote the function $f_i(x) := f(x)^i$. Note that for $z \in \Omega$, the expression
$(1+z+z^2+\cdots+z^{q-1})/q$ equals $1$ if $z=1$ and $0$ otherwise.

\begin{lemma}\label{lemma:agr} If $f: \F_q^m \mapsto \Omega$ is a function such that for some
linear function $\chi_\alpha$, $\AGR(f, \chi_\alpha) \geq {1}/{q}+\delta$. Then
$\sum_{j=1}^{q-1} \left| \widehat{f}_{j}(j \alpha) \right| \geq q\delta$.
\end{lemma}
\begin{proof} The lemma follows by noting that:
\begin{align*}
\AGR(f, \chi_\alpha)\ &=\  \Exp_x \left[  \frac{1}{q} \sum_{j=0}^{q-1} (f(x) \overline{\chi_\alpha(x)})^j   \right] \\
 &=\   \frac{1}{q} + \frac{1}{q} \sum_{j=1}^{q-1} \Exp_x \left[ f_j(x) \overline{\chi_{j \alpha}(x)}   \right] \\
 &=\  \frac{1}{q} + \frac{1}{q}  \sum_{j=1}^{q-1} \widehat{f}_j (j \alpha)\,.\qedhere
\end{align*}
\end{proof}

\subsection{The Gowers test}\label{sec:gowers-test}

For a function $f: \F_q^m \mapsto {\mathbb C}$, the Gowers Uniformity Norm $U_2$ \cite{Gowers} is defined as
 $$ \| U_2(f) \|^4 := \Exp_{x,y,z} \left[ f(x) \overline{f(x+y)}\overline{f(x+z)} f(x+y+z)   \right]\,. $$
We will study the probability that a function $f: \F_q^m \mapsto \Omega$ passes the test
$$f(x) \overline{f(x+y)}\overline{f(x+z)} f(x+y+z) =1\,,$$
for a random choice of $x,y,z$. It is thus natural to name the test as the Gowers Test.
It will be more convenient for us to think of the test equivalently
as $$\mbox{Gowers Test}: \quad\quad f(x) \overline{f(y)}\overline{f(z)} f(-x+y+z) =1\,.$$

\begin{lemma} Let $f: \F_q^m \mapsto \Omega$ be a function. Then the acceptance probability of the Gowers Test is:
$$ \Pr_{x,y,z} \left[f(x) \overline{f(y)}\overline{f(z)} f(-x+y+z) =1 \right]
   = \frac{1}{q} + \frac{1}{q} \sum_{j=1}^{q-1} \sum_\alpha  \left| \widehat{f}_j (j\alpha)\right|^4\,.   $$
\end{lemma}
\begin{proof} The acceptance probability can be expressed as
\begin{align*}
   &  \frac{1}{q} + \frac{1}{q} \sum_{j=1}^{q-1} \Exp \left[ (f(x) \overline{f(y)}\overline{f(z)} f(-x+y+z))^j   \right] \\
  =\ & \frac{1}{q} + \frac{1}{q}  \sum_{j=1}^{q-1} \sum_{\alpha, \phi, \psi, \gamma}
 \widehat{f_j}(\alpha) \overline{\widehat{f_j}(\phi)} \overline{\widehat{f_j}(\psi)}\widehat{f_j}(\gamma)
  \Exp \left[ \chi_{\alpha - \gamma}(x) \chi_{-\phi + \gamma}(y) \chi_{-\psi + \gamma}(z)    \right] \\
  =\ &  \frac{1}{q} + \frac{1}{q}  \sum_{j=1}^{q-1} \sum_\alpha \left| \widehat{f}_j (\alpha)\right|^4\,,
\end{align*}
noting that the expectation vanishes unless $\alpha = \phi = \psi = \gamma$. We can replace
$\alpha$ by $j \alpha$ without changing the summation.
\end{proof}

\subsection{The Gowers test with side conditions}

At the Inner PCP level, we need to check not only that a function $f$ is linear, \ie, $f= \chi_\alpha$ for some
$\alpha$, but also that $\alpha$ itself satisfies a given set of linear constraints. We call these
as \emph{side conditions} and modify the Gowers Test so as to incorporate these side conditions (and henceforth Gowers Test refers to one with side conditions incorporated).

\medskip
\noindent {\bf Given:}
\begin{itemize}
\item A function $f: \F_q^m \mapsto \Omega$.
\item Side conditions $h_i \cdot x = b_i, i = 1,\ldots, k$ where $h_i \in \F_q^m, \ b_i \in \F_q$. Assume that $\{h_i\}_{i=1}^k$ are linearly independent and  $H$ is their
    linear span.
\end{itemize}
\noindent {\bf The Test:}
\begin{itemize}
\item Pick $x, y, z \in \F_q^m$ at random.
\item Pick $a = (a_1,\ldots, a_k) \in \F_q^k$ at random and let $h= \sum_{i=1}^k a_i h_i$ and
      $a\cdot b := \sum_{i=1}^k a_i b_i$.
\item Accept if and only if
  $$f(x) \overline{f(y)}\overline{f(z)} f(-x+y+z+h) = \omega^{a\cdot b}\,. $$
\end{itemize}


\begin{lemma}\label{lemma:Gowers} The following hold:
\begin{enumerate}
\item The Gowers Test always passes with probability at least ${1}/{q}$.
\item If the Gowers Test passes with probability at least ${1}/{q}+\delta$, then there
exists $1 \leq j \leq q-1$ and $\alpha \in \F_q^m$ that respects the side conditions (\ie,
$\forall i, \alpha \cdot h_i = b_i$) and $| \widehat{f}_j(j \alpha) | \geq \sqrt{\delta}$.
\item If $f$ has agreement ${1}/{q}+\delta$ with a linear function $\chi_\alpha$ that respects
the side conditions (\ie,
$\forall i, \alpha \cdot h_i = b_i$), then $f$ passes the Gowers Test with probability at least
${1}/{q}+ {\delta^4}$.
\end{enumerate}
\end{lemma}
\begin{proof} The acceptance probability of the Gowers Test with side conditions is:
\begin{align*}
   &  \frac{1}{q} + \frac{1}{q} \sum_{j=1}^{q-1} \Exp \left[ (f(x) \overline{f(y)}\overline{f(z)} f(-x+y+z+h))^j   \omega^{-j a\cdot b} \right] \\
  =\ & \frac{1}{q} + \frac{1}{q}  \sum_{j=1}^{q-1} \sum_{\alpha, \phi, \psi, \gamma}
 \widehat{f_j}(\alpha) \overline{\widehat{f_j}(\phi)} \overline{\widehat{f_j}(\psi)}\widehat{f_j}(\gamma)
  \Exp \left[ \chi_{\alpha - \gamma}(x) \chi_{-\phi + \gamma}(y) \chi_{-\psi + \gamma}(z)    \right]
  \Exp_a \left[ \chi_{\gamma}(h) \omega^{-j a\cdot b}    \right]  \\
  =\ &  \frac{1}{q} + \frac{1}{q}  \sum_{j=1}^{q-1} \sum_\alpha \left| \widehat{f}_j (\alpha)\right|^4
  \cdot \Exp_a \left[ \omega^{\sum_{i=1}^k a_i (\alpha \cdot h_i - j b_i)}  \right] \\
  =\ &  \frac{1}{q} + \frac{1}{q}  \sum_{j=1}^{q-1} \ \  \sum_{\alpha| \forall i, \alpha\cdot h_i = j b_i}
  \ \ \left| \widehat{f}_j (\alpha)\right|^4 \\
   =\ &  \frac{1}{q} + \frac{1}{q}  \sum_{j=1}^{q-1} \ \  \sum_{\alpha| \forall i, \alpha\cdot h_i =  b_i}
  \ \ \left| \widehat{f}_j (j\alpha)\right|^4\,. \\
\end{align*}
The first conclusion is immediate. The second uses the fact that the sum of squared
absolute values of Fourier coefficients of $f_j$ is $1$.
 The third follows in conjunction with \expref{Lemma}{lemma:agr} and convexity of the
 function $x^4$.
\end{proof}

\subsection{The Point-Subspace test with side conditions (Inner PCP)}\label{sec:point-subspace}

\noindent {\bf Given:}
\begin{itemize}
\item A function $f: \F_q^n \mapsto \Omega$.
\item Side conditions $h_i \cdot x = b_i, i =1,\ldots, k$. Assume that $\{h_i\}_{i=1}^k$ are linearly independent and $H$ is their
    linear span.
\item A table $\left\{ {\cal T}(W) \mid  W \in {\cal C} \right\}$ where ${\cal C}$ denotes the class of $k+6$ dimensional subspaces of the form $W = D \oplus H$ for a $6$-dimensional subspace $D$ of $\F_q^n$ such that $D \cap H = \{ {\bf 0} \}$. The entry ${\cal T}(W)$ is a linear function $\chi: W \mapsto \Omega$ that respects the side conditions, \ie,
\[
\chi(x+y) = \chi(x)\cdot \chi(y)\qquad\text{and}\qquad \chi(x + h_i) =
\chi(x)\cdot  \omega^{b_i} \quad \forall x,y \in W, i=1,\ldots,k\,.
\]
\end{itemize}
\noindent {\bf The Test:}
\begin{enumerate}
\item Pick a random $W \in {\cal C}$ and let $ {\cal T}(W)$ be the linear function on $W$.
\item Pick a random $w \in W$.
\item Accept if and only if $f(w) = {\cal T}(W)(w)$.
\end{enumerate}

\subsubsection{Completeness}
Suppose $f: \F_q^n \mapsto \Omega$ is linear that respects the side
conditions, \ie, \[
f(x+y) = f(x)\cdot f(y) \qquad\text{and}\qquad f(x + h_i) = f(x)\cdot
\omega^{b_i} \quad \forall x,y \in \F_q^n, i=1,\ldots,k\,.
\]
Then letting
${\cal T}(W)$ to be the restriction $f|_W$, the point-subspace test passes with probability $1$.

\subsubsection{Soundness}

\begin{lemma} \label{lemma:inner} If the point-subspace test passes with probability
${3}/{q}$ then there exists
$1 \leq j \leq q-1$ and $\alpha\in \F_q^n$ that respects the side conditions (\ie, $\forall i,
\alpha \cdot h_i = b_i$) and $| \widehat{f}_j (j \alpha)| \geq {1}/{q^2}$.
\end{lemma}
\begin{proof} Assume that the point-subspace test passes
with probability ${1}/{q}+\delta$ where $\delta \geq {2}/{q}$. This means that:
  $$ \Exp_W \left[ \AGR\left( f|_W,  {\cal T}(W) \right) \right]  = {1}/{q}+\delta\,. $$
By \expref{Lemma}{lemma:Gowers}(1,3) and using convexity of the function $x^4$,
we conclude that
\begin{equation}
\Exp_W \left[\Pr \left[ f|_W \ \mbox{passes Gowers Test} \right] \right] \geq {1}/{q}+\delta^4.
\end{equation}
Now let us consider the probability that $f$ passes the Gowers Test. The Gowers Test
picks three points $x, y, z$ independently (from the global space $\F_q^n$). Let
$\bot$ be the event that  ${\rm span}(x,y,z,H)$ has dimension $k+3$ and
let
$\Gamma$ be the set of all such triples $(x,y,z)$. Conditional on $\bot$ happening,
the Gowers Test picks a triple in $\Gamma$ uniformly at random. An alternate way of
picking a triple in $\Gamma$ uniformly at random is to pick a subspace $W \in {\cal C}$
and then pick $(x,y,z) \in W^3$
conditional on the event that ${\rm span}(x,y,z,H)$ has dimension $k+3$. Let $\bot(W)$
be the event that ${\rm span}(x,y,z,H)$ has dimension $k+3$
when $x,y,z \in W$ are picked at random. Thus:
\begin{align*}
\Pr \left[ \text{$f$ passes Gowers Test} \mid \bot \right]\ &
 =  \ \Exp_W \bigl[ \Pr\left[ f|_W \ \mbox{passes Gowers Test} \mid \bot(W) \right]   \bigr]     \\[1ex]
 & \geq\   \Exp_W \bigl[ \Pr \left[ f|_W \ \mbox{passes Gowers Test}\right] - \Pr[\neg \bot(W)] \bigr] \\
 & \geq\   \Exp_W \bigl[ \Pr \left[ f|_W \ \mbox{passes Gowers Test}\right] \bigr]  - \frac{3}{q^4} \\
 & \geq\   \frac{1}{q} +\delta^4 -\frac{3}{q^4} \ \geq \ \frac{1}{q}+\frac{2}{q^4}\,,
 \end{align*}
where we used $\Pr[ \neg \bot(W)] \leq {3}/{q^4}$ and $\delta \geq {2}/{q}$.
Also, noting that 
\[
\Pr[\bot] \geq 1-\frac{3}{q^{n-k-2}}\,,
\]
we get that
$$\Pr\left[f \ \mbox{passes Gowers Test}\right] \geq \frac{1}{q}+\frac{1}{q^4}\,.$$
It follows now from \expref{Lemma}{lemma:Gowers}(2)
that there exists $1 \leq j \leq q-1$ and $\alpha$ such that $\forall i,
\alpha \cdot h_i = b_i$ and $| \widehat{f}_j (j \alpha)| \geq {1}/{q^2}$.
\end{proof}


\section{The composed PCP}\label{sec:composed}
We now describe the composed PCP and prove \expref{Theorem}{thm:main}.
The Outer PCP is constructed from a 3LIN instance $(X, \Phi)$
as described in \expref{Section}{sec:outer}. The 3LIN instance is either $(1-\eta)$-satisfiable
or at most $(1-\eps)$-satisfiable as per \expref{Theorem}{thm:Has97} or \expref{Theorem}{thm:kp}.
The various parameters are chosen as in \expref{Section}{sec:param}.


The verifier in the composed PCP expects the proof
to contain, for every question $V$ to the first prover, two tables $\calL_V$
and ${\cal T}_V$.

\paragraph{The tables $\calL_V$}
The table $\calL_V$ gives
the Hadamard code of the
assignment to $V$. Concretely, let $\{x_1,\ldots,x_{3k}\}$ be the variables in $V$
and $\sigma : X \mapsto \F_q$ be the global assignment that is supposed to be an
almost satisfying assignment to $(X,\Phi)$. The verifier expects, for the question $V$,
the table of values of the linear function $\calL_V: \F_q^V = \F_q^{3k}
\mapsto \Omega$ where
 $$ \calL_V(y) = \omega^{\sum_{i=1}^{3k} y_i \sigma(x_i)}\,. $$
Note that for every question $U$ to the second prover, a table $\calL_U$ that gives the
Hadamard code
of the assignment to $U$ may be expected as well. However if $U \subseteq V$, then the table $\calL_U$
is contained in the table $\calL_V$. Specifically, for any $z \in \F_q^U$,
let $z^{\uparrow}$ denote the vector obtained by extending $z$
to a vector in $\F_q^V$ be inserting
zeroes at the positions in $V\setminus U$. Then
 $$ \calL_U (z) = \omega^{\sum_{i: x_i \in U} z_i \sigma(x_i)} = \omega^{\sum_{i=1}^{3k}
    z^{\uparrow}_i \sigma(x_i)}
  = \calL_V (z^{\uparrow})\,. $$
Thus there is no need to have separate $\calL_U$ tables; we do however
think of these as virtual tables. Whenever
there are questions $V, V'$ to the first prover such that $U \subseteq V \cap V'$,
the table $\calL_U$ is contained in both the tables $\calL_V$ and $\calL_{V'}$.
We identify the
locations in these two tables that correspond to the same location in the
(virtual) $\calL_U$ table.

\paragraph{The tables $\calT_V$}
Let $V$ be a question to the first prover that consists of equations (\ie,
side conditions)
$\{h_i\cdot x = b_i\}_{i=1}^k$, the $i^{th}$ equation depending only on the
variables $(x_{3(i-1)+1}, x_{3(i-1)+2}, x_{3i})$. Except with a negligible
probability (that we may ignore) over
the choice of $V$, the $k$ equations are disjoint (\ie, do not share variables).
Thus the vectors $h_i$ are
linearly independent and let $H$ be their linear span.
The table $\calT_V$ contains, for every $k+6$ dimensional subspace $W \subseteq
\F_q^V$ such that $H\subseteq W$, a linear function $\calT_V(W) : W \mapsto \Omega$
that respects the side conditions.

\paragraph{The PCP verifier}  The verifier picks a random question
$V$ to the first prover in the Outer PCP\@. Let $\{h_i\cdot x = b_i\}_{i=1}^k$ be the side conditions and $H$ be the linear span of vectors $\{h_i\}_{i=1}^k$. Let
$\calL_V: \F_q^V \mapsto \Omega$ and $\calT_V$ be the
associated tables. The verifier picks a random $k+6$ dimensional
subspace $W$, $H \subseteq W \subseteq \F_q^V$, and a random point $w \in W$.
The verifier accepts if and only if
 $$  \calL_V(w) = \calT_V(W) (w)\,. $$

\subsection{Completeness}

Let $\sigma: X \mapsto \F_q$ be a global assignment that satisfies $1-\eta$ fraction
of equations. The table $\calL_V$ is the Hadamard code (\ie, linear function)
of the assignment $\sigma$ restricted
to $V$. The table $\calT_V$ gives, for every subspace $W$, the linear function
$\calT_V(W) = \calL_V|_W$.
The test may fail only when there is some equation in $V$ that is not satisfied
by $\sigma$. This happens with probability at most $\eta k$.



\subsection{Soundness}

Assume on the contrary that the test accepts with probability ${4}/{q}$. By an averaging argument,
for at least ${1}/{q}$ fraction of the questions $V$, the test accepts
with probability at least ${3}/{q}$. Fix any such \emph{good} $V$.
By the analysis of the Inner PCP, \expref{Lemma}{lemma:inner},
it follows that there exists $1 \leq j\leq q-1$ such that for
$f := \calL_V$, there is a Fourier coefficient $|\widehat{f_j}(j\alpha)| \geq {1}/{q^2}$ and
$\alpha$ respects the side conditions.
Note that for the uniform distribution ${\cal D}$ on $\F_q^V$,
 $$ \widehat{f_j}(j\alpha) = \Exp_{x \in {\cal D}} \left[ f_j(x)  \overline{\chi_{j\alpha}(x)}   \right]\,. $$
By the sub-code covering property, \expref{Section}{sec:param}, it follows that for some error parameter
$e, |e| \leq {1}/{(C q^2)}$,
\begin{align*}
\widehat{f_j}(j\alpha)\ &= \ \Exp_{x \in {\cal D'}} \left[  f_j(x) \overline{\chi_{j\alpha}(x)}   \right] \  + e \\
 & = \  \Exp_U \left[   \Exp_{y \in \F_q^U} \left[ f_j(y^\uparrow)
 \overline{\chi_{j\alpha}(y^\uparrow)}  \right]  \ \right] \ + e \\
 & = \  \Exp_U \left[   \Exp_{y \in \F_q^U} \left[ f_j|_U(y)
 \overline{\chi_{j\alpha_\downarrow}(y)}  \right]  \ \right] \ + e \\
 & = \  \Exp_U \left[ \widehat{f_j|_U} (j \alpha_\downarrow) \right] \ + e\,,
\end{align*}
where
$f_j|_U$ denotes the restriction of $f_j$ to $\F_q^U \subseteq \F_q^V$
and $\alpha_\downarrow$
denotes the vector obtaining by dropping coordinates of $\alpha$ in $V \setminus U$.
It follows that
 $$  \Exp_U \left[ \ \left| \widehat{f_j|_U} (j \alpha_\downarrow) \right| \right]
\geq |\widehat{f_j}(j\alpha)| - |e| \geq \frac{1}{q^2}- \frac{1}{Cq^2} \geq \frac{1}{2q^2}\,. $$
Thus, with probability at least
${1}/{(4q^2)}$ over the choice of $U$ (for the fixed good $V$; call such $(V,U)$  a \emph{good} pair),
\[
\Bigl|\widehat{f_j|_U}({j \alpha_\downarrow})\Bigr| \geq \frac{1}{4q^2}\,.
\]
Note also that
$f|_U = g = \calL_U$ which is the supposed (virtual) Hadamard code for an assignment to $U$ and
$f_j|_U = g_j$.

%
Now we derive strategies for the provers in the Outer PCP as follows: the first prover,
on receiving a question $V$, considers the function $f = \calL_V$,
lists all $\alpha$ such that for some $1 \leq j \leq q-1$, $|\widehat{f_j}(j\alpha)| \geq {1}/{q^2}$ and $\alpha$ respects the side conditions, and
then outputs a random element from the list. The list size is bounded by $q^5$.
The second prover, on receiving a question $U$, considers the function
$g = \calL_U$,
lists all $\gamma$
such that for some $1 \leq j \leq q-1$, $|\widehat{g_j}(j\gamma)| \geq {1}/{(4q^2)}$, and
outputs a random element from the list. The list size is bounded by $16q^5$.
The analysis above shows that when $(V,U)$ is a good pair, there are $\alpha$
and $\alpha_\downarrow = \gamma$ in the lists for $V$ and $U$ respectively and $\alpha$
respects the side conditions. This and the bound on the list sizes
shows that with probability at least 
\[
\frac{1}{q}\cdot \frac{1}{4q^2} \cdot\frac{1}{q^5}
\cdot \frac{1}{16q^5}\,,
\]
the provers succeed. This is a contradiction since the
soundness of the Outer PCP is at most ${1}/{q^C}$ for a large constant $C$ as in
\expref{Section}{sec:param}.
